\section{Conclusion}\label{sec:conclusion}
We introduced a splitting-off operation in (weighted) hypergraphs. Our contribution on this front is conceptualizing an appropriate notion of splitting-off in hypergraphs. 
Next, we proved that for every hypergraph there exists a local-connectivity preserving complete h-splitting-off at an arbitrary vertex from the hypergraph.  
%with polynomial number of additional hyperedges 
%and gave a strongly polynomial-time algorithm to find such a complete splitting-off at a vertex from a given \emph{weighted} hypergraph. 
Although our proof of existence follows from previously known existential results of Bern\'{a}th and Kir\'{a}ly \cite{Bernath-Kiraly} for an abstract function cover problem, our main contribution is identifying that our newly introduced notion of hypergraph splitting-off falls within the framework of their abstract function cover problem. 
Next, we designed a strongly polynomial-time algorithm to find a local-connectivity preserving complete splitting-off at a vertex. This involved substantial technical challenges to overcome. In particular, our main contribution towards the strongly polynomial-time algorithm is a strengthening of the above-mentioned existential result: we showed that there exists a local connectivity preserving complete h-splitting-off at an arbitrary vertex from the hypergraph 
\emph{in which the number of additional hyperedges is polynomial in the number of vertices}. 
%\emph{with polynomial number of additional hyperedges}. 
The strongly polynomial-time algorithm follows from our techniques to achieve this stronger existential result via standard algorithmic tools in submodularity. 
Finally, we illustrated the usefulness of the existence of local connectivity preserving complete h-splitting-off at an arbitrary vertex from a hypergraph by presenting two applications. 
Our first application is to give a constructive characterization of $k$-hyperedge-connected hypergraphs. Our second application is to give an alternative proof of an approximate min-max relation for 
the max Steiner rooted-connected orientation problem in graphs and hypergraphs. 
Two notable features of our proof of the approximate min-max relation for \emph{graphs} are that (1) it avoids the strong orientation theorem of Nash-Williams and (2) it proves a result for graphs using tools developed for hypergraphs. %We believe that these features suggest the strong 
%Our proof of existence and our algorithm to compute a local connectivity preserving complete splitting-off at an arbitrary vertex from the hypergraph are based on an algorithmic proof for the more general problem of converting a weak cover of a symmetric skew-supermodular function to a strong cover of the same function by merging disjoint hyperedges. 

Local and global connectivity preserving complete splitting-off at a vertex from a graph is a powerful operation for graphs. It finds applications in a variety of graph problems including graph orientation \cite{Mad78, Frank93:Applications-of-submod-functions, FK02, Kiraly-Lau}, connectivity augmentation \cite{Fra92, Fra94, JFJ95, Frank-book, nagamochi_ibaraki_2008_book}, 
minimum cuts enumeration \cite{NNI97, HW96, Goe01}, network design \cite{Lov76, Jor03, CS08, CFLY11}, tree packing \cite{BHTP03, Kri03, JMS03, Lau07}, congruency-constrained cuts \cite{NZ20}, and approximation algorithms for TSP \cite{GB93, BN23}. 
%These far-reaching applications have led to sustained interests in designing fast algorithms for local connectivity preserving splitting-off in graphs \cite{LY13, CLP22-soda, CLP22-stoc, CHLP23}. 
%\cite{Mad78, Frank93:Applications-of-submod-functions, FK02, Kiraly-Lau, Fra92, Fra94, JFJ95, Frank-book, nagamochi_ibaraki_2008_book, NNI97, HW96, Goe01, Jor03, CFLY11, BHTP03, Lau07, NZ20, GB93, BN23}
We believe that our local connectivity preserving complete splitting-off results for hypergraphs is likely to find future applications akin to its counterpart in graphs. 
We mention some of the open questions raised by our work:
\begin{enumerate}
    \item Our work focused on local connectivity preserving complete \emph{h-splitting-off} at a vertex from a hypergraph. We gave an example showing that local/global connectivity preserving complete \emph{g-splitting-off} at a vertex from a hypergraph may not exist (Figure \ref{figure:hypergraph-splitting-off-example-2}). 
    Are there sufficient conditions to guarantee local/global connectivity preserving complete \emph{g-splitting-off} at a vertex from a hypergraph? We recall that Lov\'{a}sz's \cite{Lov74, Lovasz-problems-book} and Mader's \cite{Mad78} results give sufficient conditions to guarantee local and global connectivity preserving complete g-splitting-off at a vertex from a \emph{graph}. 
    \iffalse
    \knote{Include these two in journal version}
    \item The \emph{rank} of a hypergraph is the size of its largest hyperedge. Hypergraphs of rank $2$ are graphs. 
    We note that complete g-splitting-off at a vertex from a graph is necessarily a graph. However, for $r\ge 3$, complete g-splitting-off/h-splitting-off at a vertex from an $r$-rank hypergraph is not necessarily an $r$-rank hypergraph---our splitting-off disjoint hyperedges operations could increase the rank. Is there a notion of splitting-off that preserves rank? Are there sufficient conditions to guarantee splitting-off operations that preserve local/global connectivity as well as rank? Such a result may help in giving a constructive characterization of $r$-rank $k$-hyperedge-connected hypergraphs. 
    \item Mader \cite{Mad82, Mad83} introduced a notion of splitting-off in directed graphs and showed that there exists global/Steiner connectivity preserving splitting-off at a vertex from a directed graph under mild conditions. This result has far reaching implications---e.g., it can be used to prove Edmonds' arborescence packing theorem \cite{Frank-book} and also finds applications in approximation algorithms for certain directed orienteering problems \cite{NR11}. Is there a generalization/counterpart of Mader's splitting-off result for directed graphs to directed hypergraphs?
    \fi
    \item As one of the applications of our splitting-off result, we presented an alternative proof of an approximate min-max relation for the max Steiner \emph{rooted}-connected orientation problem in hypergraphs. The computational complexity of a closely related hypergraph orientation problem is open: In max Steiner connected orientation problem in hypergraphs, the input is a hypergraph $G=(V, E)$ and a subset $T$ of terminals. The goal is to find the maximum $k$ and an orientation $\overrightarrow{G}$ of $G$ such that $\overrightarrow{G}$ contains $k$ hyperarc-disjoint paths from $u$ to $v$ for every pair of distinct terminals $u, v\in T$. Max Steiner connected orientation problem in \emph{graphs} is solvable in polynomial time via the Nash-Williams' strong orientation theorem. 
    Is max Steiner connected orientation problem in \emph{hypergraphs} solvable in polynomial time? 
    %Is this problem solvable in polynomial time in hypergraphs?
    %Does this problem admit a constant factor approximation in hypergraphs? 
    %\item \shubhang{Splitting off in directed hypergraphs? Such a result would help in constructive characterizations of $k$-arc-connected directed hypergraphs and $(k, \ell)$-partition-connected hypergraphs}
\end{enumerate} 

\iffalse
\knote{More questions:
\begin{enumerate}
\item Faster run-time for hypergraph splitting-off?
\item Can we use hypergraph splitting off to design an alternate algorithm to enumerate all min-cut-sets in a hypergraph? Theorem \ref{thm:splitting-off-or-trim-to-preserve-connectivity} is perhaps useful for this purpose? 
\item The \textsc{Max Steiner Rooted-Connected Orientation} Problem is $2$-approximable. What is the inapproximability? Conjecture is that it should be $(2-\epsilon)$-inapproximable. Intuition: the $2$-factor gap between connectivity and weak-partition-connectivity.
\end{enumerate}
}

\shubhang{Assuming polynomially bounded hyperedge weights, the implementation of \functionMaximizationOracleStrongCover{p} can be improved to $\Tilde{O}(|V|)$ max flow computations by computing $|V|$ Gomory-Hu trees, each using $\Tilde{O}(1)$ max-flow computations via Abboud et. al.. Perhaps a single Gomory-Hu tree of a modified hypergraph suffices? If so, then we can compute the required Gomory-Hu tree (without needing the polynomially bounded hyperedge-weights assumption) using the classical $O(|V|)$ max-flows algorithm.}
\fi


